\subsection{Primos}

\subsubsection{Crivo de Eratóstenes}

Acha todos os primos de $1$ até $n$ usando $O(n)$ memória e $O(n \log \log n)$ tempo

\begin{lstlisting}[language=c++, breaklines=true]
int n;
vector<bool> is_prime(n+1, true);
is_prime[0] = is_prime[1] = false;
for (int i = 2; i * i <= n; i++) {
    if (is_prime[i]) {
        for (int j = i * i; j <= n; j += i)
            is_prime[j] = false;
    }
}
\end{lstlisting}

É possível fazer o crivo usando apenas $O(\sqrt{n}+S)$ de memória utilizando divisão em blocos de tamanho $S$. Use $10^4 \leq S \leq 10^5$

\begin{lstlisting}[language=c++, breaklines=true]
int count_primes(int n) {
    //conta o numero de primos de 1 ateh n
    const int S = 10000;

    vector<int> primes;
    int nsqrt = sqrt(n);
    vector<char> is_prime(nsqrt + 2, true);
    for (int i = 2; i <= nsqrt; i++) {
        if (is_prime[i]) {
            primes.push_back(i);
            for (int j = i * i; j <= nsqrt; j += i)
                is_prime[j] = false;
        }
    }

    int result = 0;
    vector<char> block(S);
    for (int k = 0; k * S <= n; k++) {
        fill(block.begin(), block.end(), true);
        int start = k * S;
        for (int p : primes) {
            int start_idx = (start + p - 1) / p;
            int j = max(start_idx, p) * p - start;
            for (; j < S; j += p)
                block[j] = false;
        }
        if (k == 0)
            block[0] = block[1] = false;
        for (int i = 0; i < S && start + i <= n; i++) {
            if (block[i])
                result++;
        }
    }
    return result;
}
\end{lstlisting}

Como encontrar os números num intervalo pequeno de tamanho $10^7$ $[L, R]$ mas com $R \leq 10^{12}$? Encontre os primos até $\sqrt{R}$ e marque todos os compostos no segmento $[L, R]$. Isso roda em $O((R-L+1) \log \log(R) + \sqrt{R} \log \log \sqrt{R})$.

\begin{lstlisting}[language=c++, breaklines=true]
vector<char> segmentedSieve(long long L, long long R) {
    // generate all primes up to sqrt(R)
    long long lim = sqrt(R);
    vector<char> mark(lim + 1, false);
    vector<long long> primes;
    for (long long i = 2; i <= lim; ++i) {
        if (!mark[i]) {
            primes.emplace_back(i);
            for (long long j = i * i; j <= lim; j += i)
                mark[j] = true;
        }
    }

    vector<char> isPrime(R - L + 1, true);
    for (long long i : primes)
        for (long long j = max(i * i, (L + i - 1) / i * i); j <= R; j += i)
            isPrime[j - L] = false;
    if (L == 1)
        isPrime[0] = false;
    return isPrime;
}
\end{lstlisting}

Na verdade uma versão disso em $O((R-L+1)\log(R) + \sqrt{R})$ não exige a pre-geração dos primos:

\begin{lstlisting}[language=c++, breaklines=true]
vector<char> segmentedSieveNoPreGen(long long L, long long R) {
    vector<char> isPrime(R - L + 1, true);
    long long lim = sqrt(R);
    for (long long i = 2; i <= lim; ++i)
        for (long long j = max(i * i, (L + i - 1) / i * i); j <= R; j += i)
            isPrime[j - L] = false;
    if (L == 1)
        isPrime[0] = false;
    return isPrime;
}
\end{lstlisting}

Outra versão do crivo é em $O(n)$. O problema é que usa $32 \times$ a memória (então só faz sentido pra $n \leq 10^7$). E também não é tão "mais rápido" que o crivo normal. É até mais lento que o crivo em blocos.

O benefício de usar ele é que com ele dá pra calcular facilmente a fatoração de um número, pois a array $lp[i]$ calcula o menor fator primo de todos os números no range $[2; n]$.

\begin{lstlisting}[language=c++, breaklines=true]
const int N = 10000000;
vector<int> lp(N+1);
vector<int> pr;

for (int i=2; i <= N; ++i) {
    if (lp[i] == 0) {
        lp[i] = i;
        pr.push_back(i);
    }
    for (int j = 0; i * pr[j] <= N; ++j) {
        lp[i * pr[j]] = pr[j];
        if (pr[j] == lp[i]) {
            break;
        }
    }
}
\end{lstlisting}

Crivo mais poderoso que já vi. Ele funcionou num problema que pedia a soma dos divisores de números $\leq 10^{16}$ (olhar como resolvi na seção de funções multiplicativas) que precisou dos primos de $1$ até $10^8$.

\begin{lstlisting}[language=c++, breaklines=true]
#define MAX 100000000
#define LMT 10000

int flag[MAX / 64];
int primes[5761460], total;

#define chkC(n) (flag[n >> 6] & (1 << ((n >> 1) & 31)))
#define setC(n) (flag[n >> 6] |= (1 << ((n >> 1) & 31)))

void sieve() {
    int i, j, k;
    flag[0] |= 0;
    for (i = 3; i < LMT; i += 2)
        if (!chkC(i))
            for (j = i * i, k = i << 1; j < MAX; j += k)
                setC(j);
    primes[(j = 0)++] = 2;
    for (i = 3; i < MAX; i += 2)
        if (!chkC(i)) primes[j++] = i;
    total = j;
}
\end{lstlisting}

\subsubsection{Testes de primalidade}

Clássico $O(\sqrt{n})$

\begin{lstlisting}[language=c++, breaklines=true]
bool isPrime(int x) {
    for (int d = 2; d * d <= x; d++) {
        if (x % d == 0)
            return false;
    }
    return x >= 2;
}
\end{lstlisting}

\textbf{MillerRabin:} o deterministico passa tranquilo pra $t \leq 500$ casos e $n \leq 2^{63} - 1$

\begin{lstlisting}[language=c++, breaklines=true]
using u64 = uint64_t;
using u128 = __uint128_t;

u64 binpower(u64 base, u64 e, u64 mod) {
    u64 result = 1;
    base %= mod;
    while (e) {
        if (e & 1)
            result = (u128)result * base % mod;
        base = (u128)base * base % mod;
        e >>= 1;
    }
    return result;
}

bool check_composite(u64 n, u64 a, u64 d, int s) {
    u64 x = binpower(a, d, n);
    if (x == 1 || x == n - 1)
        return false;
    for (int r = 1; r < s; r++) {
        x = (u128)x * x % n;
        if (x == n - 1)
            return false;
    }
    return true;
};

// NAO DETERMINISTICO
bool MillerRabin(u64 n, int iter=5) { // returns true if n is probably prime, else returns false.
    if (n < 4)
        return n == 2 || n == 3;

    int s = 0;
    u64 d = n - 1;
    while ((d & 1) == 0) {
        d >>= 1;
        s++;
    }

    for (int i = 0; i < iter; i++) {
        int a = 2 + rand() % (n - 3);
        if (check_composite(n, a, d, s))
            return false;
    }
    return true;
}

// DETERMINISTICO
bool MillerRabin(u64 n) { // returns true if n is prime, else returns false.
    if (n < 2)
        return false;

    int r = 0;
    u64 d = n - 1;
    while ((d & 1) == 0) {
        d >>= 1;
        r++;
    }

    for (int a : {2, 3, 5, 7, 11, 13, 17, 19, 23, 29, 31, 37}) {
        if (n == a)
            return true;
        if (check_composite(n, a, d, r))
            return false;
    }
    return true;
}
\end{lstlisting}

\subsubsection{Fatoração}

Para $n \leq 10^{12}$ dá pra usar um crivo até $10^6$ e gerar uma lista de primos e ir testando.

Para $n \leq 10^{18}$, usar pollard

\begin{lstlisting}[language=c++, breaklines=true]
typedef long long unsigned int llui;
typedef long long int lli;
typedef long double float64;

llui mul_mod(llui a, llui b, llui m) {
    llui y = (llui)((float64)a * (float64)b / m + (float64)1 / 2);
    y = y * m;
    llui x = a * b;
    llui r = x - y;
    if ((lli)r < 0) {
        r = r + m;
        y = y - 1;
    }
    return r;
}

llui C, a, b;
llui gcd() {
    llui c;
    if (a > b) {
        c = a;
        a = b;
        b = c;
    }
    while (1) {
        if (a == 1LL) return 1LL;
        if (a == 0 || a == b) return b;
        c = a;
        a = b % a;
        b = c;
    }
}

llui f(llui a, llui b) {
    llui tmp;
    tmp = mul_mod(a, a, b);
    tmp += C;
    tmp %= b;
    return tmp;
}

llui pollard(llui n) {
    if (!(n & 1)) return 2;
    C = 0;
    llui iteracoes = 0;
    while (iteracoes <= 1000) {
        llui x, y, d;
        x = y = 2;
        d = 1;
        while (d == 1) {
            x = f(x, n);
            y = f(f(y, n), n);
            llui m = (x > y) ? (x - y) : (y - x);
            a = m;
            b = n;
            d = gcd();
        }
        if (d != n) return d;
        iteracoes++;
        C = rand();
    }
    return 1;
}

llui pot(llui a, llui b, llui c) {
    if (b == 0) return 1;
    if (b == 1) return a % c;
    llui resp = pot(a, b >> 1, c);
    resp = mul_mod(resp, resp, c);
    if (b & 1) resp = mul_mod(resp, a, c);
    return resp;
}

// Rabin-Miller primality testing algorithm
bool isPrime(llui n) {
    llui d = n - 1;
    llui s = 0;
    if (n <= 3 || n == 5) return true;
    if (!(n & 1)) return false;
    while (!(d & 1)) {
        s++;
        d >>= 1;
    }
    for (llui i = 0; i < 32; i++) {
        llui a = rand();
        a <<= 32;
        a += rand();
        a %= (n - 3);
        a += 2;
        llui x = pot(a, d, n);
        if (x == 1 || x == n - 1) continue;
        for (llui j = 1; j <= s - 1; j++) {
            x = mul_mod(x, x, n);
            if (x == 1) return false;
            if (x == n - 1) break;
        }
        if (x != n - 1) return false;
    }
    return true;
}
map<llui, int> factors;
// Precondition: factors is an empty map, n is a positive integer
// Postcondition: factors[p] is the exponent of p in prime factorization of n
void fact(llui n) {
    if (!isPrime(n)) {
        llui fac = pollard(n);
        fact(n / fac);
        fact(fac);
    } else {
        map<llui, int>::iterator it;
        it = factors.find(n);
        if (it != factors.end()) {
            (*it).second++;
        } else {
            factors[n] = 1;
        }
    }
}
\end{lstlisting}